\section{Related Work}

\subsection{KV Cache Compression for Transformers}

The quadratic memory growth of transformer KV caches has motivated extensive work on compression methods. \textbf{H2O} \cite{zhang2023h2o} introduces heavy-hitter oracle eviction, tracking accumulated attention scores across layers to identify tokens that contribute most to attention computations. By retaining a small fraction (20\%) of heavy-hitter tokens plus recent tokens, H2O achieves 29$\times$ throughput improvement on OPT-6.7B with minimal accuracy loss. However, H2O's uniform attention-based policy treats all tokens identically regardless of structural properties (query boundaries, retrieval provenance).

\textbf{StreamingLLM} \cite{xiao2024streamingllm} preserves initial attention sink tokens plus a sliding window of recent context, enabling infinite-length streaming with constant memory. While effective for conversational systems, StreamingLLM's fixed window discards distant but semantically relevant retrieved passages—problematic for multi-hop QA requiring evidence synthesis across documents.

\textbf{DynamicKV} \cite{liu2024dynamickv} extends H2O with per-layer budget allocation, observing that different transformer layers exhibit distinct attention patterns (early layers focus on local syntax, late layers on global semantics). Per-layer eviction improves efficiency by 2-3\% over global budgets. Our work is complementary—ProvenanceCache could apply per-layer budgets to provenance-aware tiers for additional gains.

A key limitation shared by these methods: they operate on attention scores computed from the original context, requiring expensive prefill passes before eviction. ProvenanceCache uses retrieval metadata available at preprocessing time, enabling eviction \textbf{before} the first generation token—eliminating redundant computation for passages destined for eviction.

\subsection{Dense Retrieval for RAG Systems}

Dense retrieval methods encode queries and passages into shared embedding spaces, ranking by semantic similarity rather than lexical overlap. \textbf{DPR} \cite{karpukhin2020dpr} trains dual BERT encoders on question-passage pairs, achieving state-of-the-art open-domain QA results by retrieving relevant contexts for reader models. \textbf{Contriever} \cite{izacard2022contriever} extends DPR to unsupervised learning via contrastive pretraining, eliminating dependence on labeled retrieval pairs.

These retrieval systems produce relevance scores (dot-product similarities) and passage boundaries as byproducts of retrieval. However, downstream generation models (reader LLMs) typically discard this metadata after concatenating retrieved passages into a single context string. Our work demonstrates that retrieval scores \textbf{generalize} from passage ranking (retrieval task) to attention prediction (generation task), with Contriever showing $\rho$=0.612 Spearman correlation—57\% stronger than lexical BM25 ($\rho$=0.391).

\subsection{Diversity in Information Retrieval}

Diversity-aware ranking prevents redundant results in search systems. \textbf{Maximal Marginal Relevance (MMR)} (Carbonell \& Goldstein, 1998) balances relevance and novelty via a weighted combination: $\lambda \times$ relevance + $(1-\lambda) \times$ diversity (embedding cosine distance from already-selected documents). MMR with $\lambda$=0.5 produces diverse result sets without excessive relevance sacrifice.

Coverage-based ranking \cite{clarke2008novelty} models query intent as a distribution over subtopics, rewarding documents that cover under-represented aspects. These methods target retrieval ranking (top-k selection from large candidate pools). ProvenanceCache extends diversity principles to \textbf{cache eviction}—memory-constrained retention where all passages are initially in cache, and we must select which to preserve.

Our work shows diversity matters asymmetrically by task type: +14.71\% F1 gain on multi-hop QA (bridging distant facts across dissimilar passages) versus +6.16\% on single-hop factoid QA (top-k relevance ranking sufficient). This finding validates that multi-hop reasoning is a \textbf{coverage problem} (maximize diverse evidence span), not a \textbf{ranking problem} (maximize top-k scores).

\subsection{Multi-Hop Question Answering}

Multi-hop QA requires synthesizing evidence across multiple documents. \textbf{HotpotQA} \cite{yang2018hotpotqa} introduced the bridge entity task: answering ``What nationality is the director of film X?'' requires retrieving (1) ``X was directed by Y'' and (2) ``Y has nationality Z.'' Both passages are semantically relevant to the query, but neither alone suffices—diversity preserves both rather than overweighting the first high-relevance hit.

\textbf{MuSiQue} \cite{trivedi2022musique} extends multi-hop reasoning to disconnected reasoning chains requiring 2-4 hops across compositional questions. These datasets motivate our diversity-aware eviction policy: naïve top-k retention by relevance over-allocates cache budget to redundant passages about Entity A, evicting critical bridging facts about Entity B.

Our experiments on LongBench multi-doc QA (subset of HotpotQA + narrativeqa + triviaqa) demonstrate that MMR diversity scoring prevents this failure mode, achieving 2.4$\times$ larger gains on multi-hop questions than single-hop factoid QA.

\subsection{Relationship to Prior Work}

ProvenanceCache differs from prior KV cache compression in three ways:

\begin{enumerate}
\item \textbf{Metadata source}: Uses retrieval provenance (passage boundaries, Contriever scores) rather than runtime attention (H2O, DynamicKV) or positional heuristics (StreamingLLM).

\item \textbf{Eviction timing}: Applies provenance-based eviction \textbf{before} first generation token (at preprocessing), not after accumulated attention tracking.

\item \textbf{Task-aware diversity}: Incorporates MMR diversity metric to prevent redundant high-relevance passage retention on multi-hop tasks.
\end{enumerate}

To our knowledge, this is the first work to condition KV cache eviction on retrieval metadata and empirically validate that retrieval scores predict attention patterns during generation ($\rho$=0.612 correlation).
